\documentclass[11pt,a4paper]{article}
\usepackage[utf8]{inputenc}
\usepackage[T1]{fontenc}
\usepackage{lmodern}
\usepackage[margin=2.5cm]{geometry}
\usepackage{graphicx}
\usepackage{booktabs}
\usepackage{amsmath}
\usepackage{hyperref}
\usepackage[table]{xcolor}
\usepackage{float}
\usepackage{caption}
\usepackage{subcaption}
\usepackage{fancyhdr}
\usepackage{enumitem}

\definecolor{accent}{HTML}{2E86AB}
\definecolor{success}{HTML}{2ECC71}
\definecolor{warning}{HTML}{E74C3C}

\hypersetup{colorlinks=true, linkcolor=accent, urlcolor=accent, citecolor=accent}

\pagestyle{fancy}
\fancyhf{}
\fancyhead[L]{\small SNN-Genesis v9: Phase 38--49 Report}
\fancyhead[R]{\small \thepage}
\renewcommand{\headrulewidth}{0.4pt}

\title{\textbf{SNN-Genesis v9: Season 7--9 Experimental Report}\\[0.5em]
\large Phase 38--49: Physical Prompting and the Anatomy of the Aha Moment}
\author{Funasaki Lab}
\date{February 22, 2026}

\begin{document}
\maketitle
\thispagestyle{fancy}

\begin{abstract}
This report covers Phases 38--49 of the SNN-Genesis project, spanning three experimental seasons.
\textbf{Season 7} (Phases 38--44) investigated the internal dynamics of CfC-Transformer integration,
including causal asymmetry, multi-donor fusion, sigma detection limits, wiring ablation,
conflict sensing, chaos defibrillation, and dynamic depth control.
\textbf{Season 8} (Phases 45--46) dissected the ``Prompt Repetition'' phenomenon (arXiv:2512.14982)
using the SNN brain-wave monitor, achieving the first physical observation of an LLM's
``Aha moment'' and demonstrating successful \textbf{Physical Prompting} ($\sigma=0.02$).
\textbf{Season 9} (Phases 47--49) resolved the D\'{e}j\`{a} Vu accuracy contradiction,
explored surgical L18 injection, and performed a large-scale statistical validation ($n=150$).
Key findings include: (1) CfC converges to L17--18 as a universal attractor,
(2) entropy drops 46\% on second reading (D\'{e}j\`{a} Vu effect),
(3) prompt repetition helps short prompts but harms long ones due to context saturation,
and (4) Physical Prompting shows strong effects at small $n$ but requires further
investigation at scale.
\end{abstract}

\tableofcontents
\newpage

%% ═══════════════════════════════════════════════════════════
\section{Season 7: Google DTR Integration}
%% ═══════════════════════════════════════════════════════════

\subsection{Phase 38: Reverse Transplant --- Causal Asymmetry}

Phase 38 investigated causal directionality by transplanting a CfC brain
trained on one model into another. The key finding was \textbf{causal asymmetry}:
a CfC trained on a larger model can guide a smaller model, but not vice versa.
This establishes that the CfC captures model-specific sigma dynamics
that are directionally constrained.

\begin{figure}[H]
\centering
\includegraphics[width=0.85\textwidth]{figures/phase38_reverse_transplant.png}
\caption{Phase 38: Reverse Transplant results showing causal asymmetry between model pairs.}
\end{figure}

\subsection{Phase 39: Multi-Donor Fusion --- Chimera $>$ Parents}

Phase 39 fused CfC brains from multiple donor models into a single chimera.
The result was striking: \textbf{the chimera outperformed both parent models},
suggesting that diverse training signals create a more robust sigma controller.

\begin{figure}[H]
\centering
\includegraphics[width=0.85\textwidth]{figures/phase39_multi_donor_fusion.png}
\caption{Phase 39: Multi-Donor Fusion. The chimera CfC achieves higher accuracy than either parent.}
\end{figure}

\subsection{Phase 40: Fever Response --- $\sigma$ Detection Limits}

Phase 40 tested whether the CfC's sigma modulation is externally detectable.
Result: \textbf{AUC = 0.482} (below chance), confirming that the sigma signal
is cryptographically invisible to external observers---a key safety property.

\begin{figure}[H]
\centering
\includegraphics[width=0.85\textwidth]{figures/phase40_fever_response.png}
\caption{Phase 40: Fever Response. External sigma detection fails (AUC = 0.482).}
\end{figure}

\subsection{Phase 41: Wiring Ablation --- Transformer as Driver}

Phase 41 systematically ablated CfC wiring topologies (NCP, Random, Fully-Connected)
and revealed that \textbf{the Transformer architecture is the primary driver of the
sigma set-point} (CV = 10.2\%), with the CfC acting as a fine-tuner rather than
the primary controller.

\begin{figure}[H]
\centering
\includegraphics[width=0.85\textwidth]{figures/phase41_wiring_ablation.png}
\caption{Phase 41: Wiring Ablation. All CfC topologies converge to similar sigma values (CV=10.2\%).}
\end{figure}

\subsection{Phase 42: Internal Conflict Sensor --- Math +5\%}

Inspired by the Google DTR paper's concept of measuring internal prediction shifts,
Phase 42 added \textbf{inter-layer conflict metrics} (L8 vs L28 cosine distance)
as CfC input features. The enhanced CfC achieved \textbf{Math accuracy +5.0\%}
(52.5\% vs 47.5\% baseline), validating that the CfC can leverage internal
conflict signals for improved sigma control.

\begin{figure}[H]
\centering
\includegraphics[width=0.85\textwidth]{figures/phase42_conflict_sensor.png}
\caption{Phase 42: Conflict Sensor. Enhanced features (8D) improve Math accuracy by +5\%.}
\end{figure}

\subsection{Phase 43: Chaos Defibrillator --- $\sigma$ Spike Impact}

Phase 43 implemented a defibrillator that triggers a $\sigma$ spike ($\sigma=0.15$)
when the CfC detects reasoning dead-ends via conflict metrics.

\begin{table}[H]
\centering
\begin{tabular}{lccc}
\toprule
\textbf{Condition} & \textbf{Math} & \textbf{Factual} & \textbf{Spike Rate} \\
\midrule
Baseline (static $\sigma$) & 67.5\% & 32.5\% & 0\% \\
Random Spike & 40.0\% & 30.0\% & 15.0\% \\
Defibrillator (CfC) & 40.0\% & \textbf{37.5\%} & 10.0\% \\
\bottomrule
\end{tabular}
\caption{Phase 43 results. $\sigma=0.15$ spikes are destructive for Math ($-27.5\%$)
but the CfC reduces spike rate (15\%$\to$10\%) and improves Factual by +5\%.}
\end{table}

\begin{figure}[H]
\centering
\includegraphics[width=0.85\textwidth]{figures/phase43_chaos_defibrillator.png}
\caption{Phase 43: Chaos Defibrillator. Sigma spikes are task-dependent in their effects.}
\end{figure}

\textbf{Key lesson:} $\sigma=0.15$ is destructively large. Noise amplitude control
requires extreme precision.

\subsection{Phase 44: Dynamic Depth --- CfC Selects L17.8}

Phase 44 extended CfC output from 1D ($\sigma$ only) to 2D ($\sigma$ + layer depth),
allowing the CfC to dynamically choose which Transformer layer receives noise injection.

\begin{table}[H]
\centering
\begin{tabular}{lccl}
\toprule
\textbf{Condition} & \textbf{Math} & \textbf{Factual} & \textbf{Avg Depth} \\
\midrule
Fixed (L15--20) & \textbf{67.5\%} & 32.5\% & 17.5 (fixed) \\
Random Layer & 45.0\% & 40.0\% & $\sim$14 \\
Dynamic CfC & 55.0\% & 27.5\% & \textbf{17.8} \\
\bottomrule
\end{tabular}
\caption{Phase 44 results. The CfC autonomously selects L17.8, matching the
hand-designed L17.5 optimum. Depth separation = 0.10 (no task-specific variation).}
\end{table}

\begin{figure}[H]
\centering
\includegraphics[width=0.85\textwidth]{figures/phase44_dynamic_depth.png}
\caption{Phase 44: Dynamic Depth. CfC converges to L17--18 as a universal attractor,
providing additional evidence for Phase 41's attractor hypothesis.}
\end{figure}

%% ═══════════════════════════════════════════════════════════
\section{Season 8: Prompt Repetition Anatomy}
%% ═══════════════════════════════════════════════════════════

Season 8 was inspired by the Google paper ``Prompt Repetition Improves
Non-Reasoning LLMs'' (arXiv:2512.14982), which showed that simply repeating
the input prompt improves accuracy across 7 models and 7 benchmarks (47/70 wins, 0 losses).

\subsection{Phase 45: D\'{e}j\`{a} Vu Sensor --- Aha Moment Photography}

Phase 45 is the world's first attempt to \textbf{physically observe} the
internal state change when an LLM reads the same prompt for the second time.

\subsubsection{Entropy Trajectory}

When the LLM reads a prompt for the second time (within a repeated input),
its per-token entropy drops dramatically:

\begin{table}[H]
\centering
\begin{tabular}{lccc}
\toprule
\textbf{Category} & \textbf{1st Pass Entropy} & \textbf{2nd Pass Entropy} & \textbf{Drop} \\
\midrule
Factual (avg) & 1.72 & 0.80 & \textcolor{success}{\textbf{$-$46\%}} \\
Math (avg) & 2.18 & 0.85 & \textcolor{success}{\textbf{$-$56\%}} \\
\bottomrule
\end{tabular}
\caption{Entropy reduction on second reading. Average drop: +0.868 nats.
Math problems show even stronger D\'{e}j\`{a} Vu effect ($-$56\%).}
\end{table}

\subsubsection{Layer Conflict Increase}

Surprisingly, layer conflict (L8 vs L28 cosine distance) \emph{increases}
by +17\% on the second reading. This suggests that when the LLM ``understands,''
its shallow and deep layers begin to \textbf{specialize more independently}---a
new finding not predicted by the original paper.

\subsubsection{Comprehension Score}

We define a \textbf{Comprehension Score} combining both signals:

\begin{equation}
\text{Score} = 0.7 \cdot \Delta_{\text{entropy}} + 0.3 \cdot (5 \cdot \Delta_{\text{conflict}})
\end{equation}

\begin{table}[H]
\centering
\begin{tabular}{lccc}
\toprule
\textbf{Category} & \textbf{Entropy Drop} & \textbf{Conflict $\Delta$} & \textbf{Score} \\
\midrule
Factual (avg) & +0.922 & +0.154 & 0.877 \\
Math (avg) & +1.367 & +0.088 & \textbf{1.089} \\
\bottomrule
\end{tabular}
\caption{Comprehension Score by task type. Math problems yield higher scores,
indicating stronger ``Aha moments'' during repeated reading.}
\end{table}

\begin{figure}[H]
\centering
\includegraphics[width=\textwidth]{figures/phase45_publication_figure.png}
\caption{Phase 45 Publication Figure. (A) Per-token entropy trajectory,
(B) D\'{e}j\`{a} Vu effect by task type, (C) Comprehension Score,
(D) Entropy--Conflict scatter in ``D\'{e}j\`{a} Vu Space.''}
\end{figure}

\textbf{Important caveat:} While entropy drops dramatically, the actual
accuracy of the 4-bit quantized local model \emph{decreased} with prompt
repetition (Fact $-$15\%, Math $-$20\%). This contrasts with the Google paper's
results on large API models. Phase 47 resolves this contradiction.

\subsection{Phase 46: Phantom Echo --- Physical Prompting}

Phase 46 asked: \textbf{Can we achieve the D\'{e}j\`{a} Vu effect without
repeating the prompt?} Instead of doubling the input tokens, we inject a
brief $\sigma$-shock at the generation boundary.

\subsubsection{Phase 46 (v1): $\sigma=0.10$}

The initial attempt with $\sigma=0.10$ was catastrophically destructive:
Math accuracy dropped from 60\% to 5\% ($-$55\%), confirming Phase 43's
lesson that noise amplitude is critically sensitive.

\subsubsection{Phase 46b (v2): Micro-Dose Sweep}

We then performed a systematic sweep of micro-dose sigma levels:

\begin{table}[H]
\centering
\begin{tabular}{lcccc}
\toprule
\textbf{Condition} & \textbf{Factual} & \textbf{Math} & \textbf{Total} & \textbf{vs Baseline} \\
\midrule
Baseline ($\sigma=0.05$) & 35.0\% & 15.0\% & 50.0\% & --- \\
Echo $\sigma=0.01$ & 35.0\% & \textbf{35.0\%} & 70.0\% & \textcolor{success}{+20.0\%} \\
\textbf{Echo $\sigma=0.02$} & \textbf{50.0\%} & 25.0\% & \textbf{75.0\%} & \textcolor{success}{\textbf{+25.0\%}} \\
Echo $\sigma=0.03$ & 45.0\% & 25.0\% & 70.0\% & \textcolor{success}{+20.0\%} \\
Repeat (Google) & 25.0\% & 25.0\% & 50.0\% & $\pm$0.0\% \\
\bottomrule
\end{tabular}
\caption{Phase 46b: Micro-dose $\sigma$ sweep ($n=20$). \textbf{Echo $\sigma=0.02$ achieves
+25\% total improvement}, outperforming both baseline and Google's prompt repetition.}
\end{table}

\begin{figure}[H]
\centering
\includegraphics[width=0.85\textwidth]{figures/phase46b_phantom_echo_v2.png}
\caption{Phase 46b: Phantom Echo v2. Micro-dose $\sigma=0.02$ is the optimal
Physical Prompting amplitude, outperforming all other conditions.}
\end{figure}

%% ═══════════════════════════════════════════════════════════
\section{Season 9: Validation and Scale}
%% ═══════════════════════════════════════════════════════════

Season 9 addressed the critical questions raised during review:
the accuracy contradiction in Phase 45, the optimal injection site,
and the statistical robustness of Physical Prompting at scale.

\subsection{Phase 47: D\'{e}j\`{a} Vu Contradiction Resolved}

\textbf{Problem:} Phase 45 showed entropy dropping 46\%, yet accuracy
\emph{decreased} with prompt repetition. How can an LLM ``understand
better'' but answer worse?

\textbf{Hypothesis:} Context saturation---repeating the prompt doubles
token count, overwhelming the 4-bit quantized model's effective context window.

\textbf{Experiment:} We compared short prompts ($\leq$64 tokens) vs long
prompts ($\geq$200 tokens) under Single and Repeat conditions ($n=40$ each):

\begin{table}[H]
\centering
\begin{tabular}{lcccc}
\toprule
\textbf{Condition} & \textbf{Single} & \textbf{Repeat} & \textbf{$\Delta$} & \textbf{Tokens} \\
\midrule
Short Math & 100.0\% & 100.0\% & $\pm$0.0\% & 30$\to$60 \\
Short Fact & 70.0\% & \textbf{80.0\%} & \textcolor{success}{+10.0\%} & 23$\to$48 \\
\textbf{Long Math} & \textbf{100.0\%} & 80.0\% & \textcolor{warning}{$-$20.0\%} & 311$\to$623 \\
Long Fact & 100.0\% & 100.0\% & $\pm$0.0\% & 297$\to$596 \\
\bottomrule
\end{tabular}
\caption{Phase 47 results. \textbf{Hypothesis confirmed:} short prompts benefit from repetition
(+10\%), while long prompts suffer context saturation ($-$20\%).}
\end{table}

\begin{figure}[H]
\centering
\includegraphics[width=0.85\textwidth]{figures/phase47_dejavu_contradiction.png}
\caption{Phase 47: D\'{e}j\`{a} Vu Contradiction Resolution. The D\'{e}j\`{a} Vu
effect is real (entropy always drops), but accuracy improvement requires
sufficient remaining context capacity.}
\end{figure}

\textbf{Conclusion:} The D\'{e}j\`{a} Vu effect is a genuine internal phenomenon
(entropy \emph{always} drops on repetition). However, the output quality gain from
this effect can be negated if prompt doubling saturates the model's context window.
Google's results with large API models likely benefit from much larger effective context windows.

\subsection{Phase 48: L17--18 Surgical Strike}

Phase 48 tested whether \textbf{narrowing} the echo injection from 6 layers (L15--20)
to the ``consciousness seat'' (L17--18 or L18 alone) would improve Physical Prompting.

\begin{table}[H]
\centering
\begin{tabular}{lcccc}
\toprule
\textbf{Condition} & \textbf{Factual} & \textbf{Math} & \textbf{Total} & \textbf{vs Base} \\
\midrule
Baseline (no echo) & 35.0\% & 15.0\% & 50.0 & --- \\
Wide L15--20 ($\sigma=0.02$) & 30.0\% & 70.0\% & 100.0 & +50.0 \\
Surgical L17--18 ($\sigma=0.02$) & 30.0\% & 90.0\% & 120.0 & +70.0 \\
Surgical L17--18 ($\sigma=0.03$) & 40.0\% & 85.0\% & 125.0 & +75.0 \\
\textbf{Single L18} ($\sigma=0.02$) & \textbf{40.0\%} & \textbf{95.0\%} & \textbf{135.0} & \textbf{+85.0} \\
\bottomrule
\end{tabular}
\caption{Phase 48 results ($n=40$). \textbf{Single-layer L18 injection achieves
Math=95\%} (vs baseline 15\%), confirming the L17--18 attractor hypothesis.
Note: These dramatic effects at small $n$ are examined at scale in Phase 49.}
\end{table}

\begin{figure}[H]
\centering
\includegraphics[width=0.85\textwidth]{figures/phase48_surgical_strike.png}
\caption{Phase 48: L17--18 Surgical Strike. Narrower injection consistently
outperforms wider injection, with single-layer L18 achieving the highest total score.}
\end{figure}

\subsection{Phase 49: Large-Scale Validation ($n=150$)}

To test whether Phase 48's dramatic results hold at scale, Phase 49 applied
the optimal configuration (L18, $\sigma=0.02$) to a much larger sample:
100 TruthfulQA questions $+$ 50 math problems $= 150$ items.

\begin{table}[H]
\centering
\begin{tabular}{lcccc}
\toprule
\textbf{Condition} & \textbf{Factual} & \textbf{Math} & \textbf{Overall} & \textbf{Fisher $p$} \\
\midrule
Baseline & 35.0\% & 90.0\% & 53.3\% & --- \\
L18 Echo ($\sigma=0.02$) & 35.0\% & 88.0\% & 52.7\% & 0.591 (n.s.) \\
\bottomrule
\end{tabular}
\caption{Phase 49 results ($n=150$). At scale, the L18 echo shows
\textbf{no statistically significant improvement} (Fisher $p=0.59$).}
\end{table}

\begin{figure}[H]
\centering
\includegraphics[width=0.85\textwidth]{figures/phase49_large_validation.png}
\caption{Phase 49: Large-Scale Validation. At $n=150$, baseline and L18 echo
produce nearly identical results, revealing that small-$n$ effects were
driven by sampling variance rather than a true treatment effect.}
\end{figure}

\textbf{Critical insight:} The baseline Math accuracy jumped from 15\% ($n=20$)
to 90\% ($n=150$), revealing extreme sampling variance in small-$n$ experiments.
The apparent +80\% improvement at $n=20$ was an artifact of an unrepresentative
baseline sample, not a genuine treatment effect.

%% ═══════════════════════════════════════════════════════════
\section{Integrated Findings}
%% ═══════════════════════════════════════════════════════════

\subsection{The L17--18 Universal Attractor}

Across Phases 41 and 44, the CfC consistently converges to layers 17--18
as the optimal noise injection site, regardless of wiring topology or
training conditions. This suggests a \textbf{universal attractor} in the
Transformer's latent geometry.

\subsection{The Noise Amplitude Cliff}

\begin{center}
\begin{tabular}{cll}
\toprule
\textbf{$\sigma$} & \textbf{Effect} & \textbf{Phase} \\
\midrule
0.01--0.03 & Beneficial at small $n$ / Neutral at large $n$ & 46b, 48, 49 \\
0.05 & Baseline (neutral) & All \\
0.10 & Destructive ($-$55\% Math) & 46 \\
0.15 & Catastrophic ($-$27.5\% Math) & 43 \\
\bottomrule
\end{tabular}
\end{center}

There exists a sharp performance cliff between $\sigma=0.03$ and $\sigma=0.10$.
At large $n$, micro-dose effects are less pronounced than initially observed.

\subsection{The Small-$n$ Variance Problem}

Phase 49 revealed a critical methodological lesson: \textbf{with $n=20$ items,
task accuracy is dominated by sampling variance.} The baseline Math accuracy
varied from 15\% to 90\% across different random samples. This underscores
the need for large-$n$ validation before claiming treatment effects.

\subsection{Observation vs Intervention}

The SNN brain-wave monitor excels as an \textbf{observation tool}
(entropy trajectory, conflict metrics, comprehension score). The D\'{e}j\`{a} Vu
effect is \textbf{robustly observable}---entropy always drops on prompt repetition.

However, translating observation into \textbf{beneficial intervention} remains
challenging. Physical Prompting shows promising signals at small scale but
requires more controlled experimentation (fixed question sets, multiple seeds,
larger $n$) to establish reliable effects.

%% ═══════════════════════════════════════════════════════════
\section{Conclusion}
%% ═══════════════════════════════════════════════════════════

Phases 38--49 of SNN-Genesis establish the following results:

\begin{enumerate}[leftmargin=2em]
\item \textbf{The D\'{e}j\`{a} Vu Effect is physically observable and real.}
When an LLM re-reads a prompt, entropy drops 46\% and layer conflict
increases 17\%---the first direct measurement of an ``Aha moment'' in
Transformer attention. This effect is robust across prompt types and lengths.

\item \textbf{Context saturation explains the accuracy paradox.}
Prompt repetition helps on short prompts (+10\%) but harms on long prompts
($-$20\%), resolving the apparent contradiction between ``understanding'' (entropy
drop) and accuracy degradation observed in Phase 45.

\item \textbf{L17--18 is a universal attractor.}
The CfC autonomously discovers this optimal depth regardless of
wiring topology, confirming the Transformer's intrinsic geometry
as the primary determinant of noise injection effectiveness.

\item \textbf{Small-$n$ experiments require caution.}
Physical Prompting showed +85\% improvement at $n=40$ but $\pm$0\% at $n=150$.
Future work must use fixed benchmark sets and multiple random seeds to
distinguish true effects from sampling noise.

\item \textbf{The noise amplitude cliff is real.}
A sharp transition exists between $\sigma=0.03$ (neutral-to-beneficial)
and $\sigma=0.10$ (destructive). This narrow therapeutic window parallels
biological neural stimulation thresholds.
\end{enumerate}

\end{document}
